\begin{lemma}[Restriction of a Tashkinov augmentation sequence]
Let \(S,B,C\) be tree-sequences for the same multigraph and distinguished
first edge, and let \(\varphi\) be one fixed partial edge coloring.  Suppose
that \(C\) is obtained from \(S\) by a
\(\varphi\)-\noderef{TashkinovAugmentationSequence}, that \(S\) is an exact
ordered \noderef{TreeSequenceInitialSegment} of \(B\), and that \(B\) is an
exact ordered \noderef{TreeSequenceInitialSegment} of \(C\).  Then \(B\) is
obtained from \(S\) by a
\(\varphi\)-\noderef{TashkinovAugmentationSequence}.
\end{lemma}
\begin{proof}
We first record two consequences of the exact ordered initial-segment
definition.  First, if \(X\) is an exact ordered initial segment of \(Y\)
and the two tree-sequences have the same length, then \(X=Y\).  Indeed, the
lengths agree, and the vertex and edge equalities in
\noderef{TreeSequenceInitialSegment} identify every entry at the same
index.  The remaining fields of the tree-sequence structures are proofs of
the attachment and injectivity conditions, so proof irrelevance identifies
those fields as well.

Second, suppose that \(X\) and \(Y\) are exact ordered initial segments of
the same tree-sequence \(Z\), and that the length of \(Y\) is at most the
length of \(X\).  Then \(Y\) is an exact ordered initial segment of \(X\).
For each vertex or edge index of \(Y\), its entry in \(Y\) equals the entry
of \(Z\) at that index, while the entry of \(X\) at the same index also
equals that entry of \(Z\).  Thus the entries of \(Y\) and \(X\) agree
through the length of \(Y\), which is precisely
\noderef{TreeSequenceInitialSegment}.

By \noderef{TashkinovAugmentationSequence}, the given passage from \(S\) to
\(C\) is a finite reflexive-transitive chain of one-step
\noderef{TashkinovAugmentation}s, all under \(\varphi\).  We induct on this
chain while allowing \(B\) to be any exact ordered intermediate segment.
If the chain is empty, then \(C=S\).  The two assumed initial-segment length
inequalities give
\[
  |B|\le |C|=|S|\le |B|.
\]
Hence \(B\) and \(S\) have equal length, and the first consequence above,
applied to the initial-segment relation from \(B\) to \(C=S\), gives
\(B=S\).  The empty augmentation sequence therefore proves the conclusion.

For the induction step, write the final one-step augmentation as
\(A\longrightarrow C\).  By \noderef{TashkinovAugmentation}, \(A\) is an
exact ordered initial segment of \(C\) and
\[
  |C|=|A|+1.
\]
If \(|B|\le |A|\), the second consequence above, applied to the two initial
segments \(A,B\) of \(C\), shows that \(B\) is an exact ordered initial
segment of \(A\).  The induction hypothesis restricts the earlier chain
from \(S\) to \(A\) to a chain from \(S\) to \(B\), as required.

It remains that \(|B|>|A|\).  Since \(B\) is an initial segment of \(C\),
we also have \(|B|\le |C|=|A|+1\), and hence \(|B|=|C|\).  The first
consequence above now gives \(B=C\).  Appending the displayed final
\(\varphi\)-\noderef{TashkinovAugmentation} to the earlier chain from \(S\)
to \(A\) yields the required augmentation sequence from \(S\) to \(B=C\).
This completes the induction.
\end{proof}
